% ============================================================================
\section{Software Design}
\label{sec:software}
% ============================================================================

\subsection{Toolchain and project layout}
The firmware is an ESP-IDF~v6.0.3 project for the ESP32-S3 (\file{firmware/}), written in C11. It
builds with no warnings and uses about 26\,\% of the application partition. The same sources also
build with the wokwi.com online builder (ESP-IDF~v5.3), which is how the shared Wokwi link runs
(\cref{app:wokwi-link}).

\begin{lstlisting}[style=console]
firmware/
  main/              main.c (boot order), Kconfig.projbuild (tunable settings)
  components/
    common/          types, pin map, task table, SHA-256
    core_logic/      access_policy, room_logic, auth_fsm   <- pure, unit tested
    drivers/         i2c_bus, lcd2004, ds1307, mfrc522, buttons, pwm, actuators, light_sensor
    bus/             messages between terminals
    storage/         SD card gatekeeper + hash-chained log
    services/        users, config, occupancy, auth, clock (+ nvs_store helper)
    ui/              LCD, buzzer, LED task
    nodes/           hub, hub_access, room_node, door, rfid
    cli/             serial menu: cmd_monitor, cmd_admin, cmd_sim
  test/host/         PC unit tests for core_logic
  tools/             gen_diagram.py (generates diagram.json)
  diagram.json, wokwi.toml   Wokwi simulation
\end{lstlisting}

Every source file starts with a comment block that says what the module does and why, and every
public function is documented with its parameters, return value and the tasks allowed to call it.

\subsection{FreeRTOS tasks}
\label{sec:sw-tasks}
Work is divided into tasks so that a slow job (an SD write, a buzzer tone) never delays a fast
one (a door opening). Stack sizes and priorities are kept together in
\file{components/common/include/task_config.h}; a higher number is a higher priority.

\begin{center}
\small
\begin{tabularx}{\linewidth}{@{}lccL@{}}
\toprule
Task & Prio & Stack & Job \\
\midrule
\code{rfid}    & 6 & 3\,KB & Polls both RFID readers every 100\,ms, reads a newly presented card's
  UID, sends \code{MSG_CARD_SCAN}; re-initialises a reader that stopped answering. \\
\code{buttons} & 6 & 3\,KB & Debounces the break-glass button (interrupt $\rightarrow$ queue
  $\rightarrow$ task), sends \code{MSG_BREAK_GLASS}. \\
\code{hub}     & 5 & 6\,KB & Entrance terminal: access decisions, admin login, room supervision,
  energy log. \\
\code{room1}, \code{room2} & 4 & 4\,KB & Room terminals: sense every 200\,ms, decide the lamp
  level, report to the hub. \\
\code{storage} & 3 & 8\,KB & SD card gatekeeper: appends, tails and verifies the log. \\
\code{ui}      & 2 & 4\,KB & Redraws the LCD every 500\,ms, plays sounds, times the LED. \\
\code{cli}     & 1 & 8\,KB & Serial command menu (lowest priority: a human is typing). \\
\bottomrule
\end{tabularx}
\end{center}
The door auto-close uses one \code{esp_timer} per door (a callback, not a task).

Input tasks have the highest priority so no card or button press is missed. The storage, UI and CLI tasks
are below the terminals because they only consume results.

\subsection{Boot sequence}
\file{main/main.c} brings the system up in dependency order, so each step can use the ones before:
\begin{enumerate}
  \item NVS (flash storage), the I2C bus, and the clock (time from the DS1307 RTC, or the build
        time as a fallback).
  \item The storage task (mounts the SD card and resumes the log chain), then a SHA-256 self-test.
        A failure is logged as \code{ALERT code=CRYPTO_FAIL}.
  \item Shared state: users, configuration, occupancy ledger, auth.
  \item The bus and the light sensor ADC.
  \item The hub task, then room nodes 1 and 2.
  \item A \code{BOOT} log line with the reset reason, then the serial CLI.
\end{enumerate}

\subsection{The node bus}
\label{sec:sw-bus}
\file{components/bus/} is how the terminals talk. Each node (\code{HUB}, \code{ROOM1},
\code{ROOM2}) has an inbox, a FreeRTOS queue of 16 messages. A message is a tagged union:

\begin{center}
\small
\begin{tabularx}{\linewidth}{@{}llL@{}}
\toprule
Type & Direction & Payload \\
\midrule
\code{MSG_CARD_SCAN} & RFID reader $\rightarrow$ hub & card UID, side (IN/OUT) \\
\code{MSG_BREAK_GLASS} & button $\rightarrow$ hub & -- \\
\code{MSG_ROOM_STATE} & room $\rightarrow$ hub & occupied, daylight, level, lux, mode, energy used/saved \\
\code{MSG_HEARTBEAT} & room $\rightarrow$ hub & uptime \\
\code{MSG_LIGHT_CMD} & hub $\rightarrow$ room & light mode (AUTO, ON, OFF) \\
\bottomrule
\end{tabularx}
\end{center}

\code{bus_send()} stamps each message with its source and the next \textbf{sequence number} for
that sender. \code{bus_recv()} remembers the last sequence number seen from each sender and drops
any message that is not newer, so a duplicated or replayed frame is never processed twice. The
number dropped is shown by \code{status}.

\subsection{RFID reader driver}
\label{sec:sw-rfid}
\file{components/drivers/mfrc522.c} talks to the MFRC522 chip over SPI. The chip is a radio front
end: the ESP32 puts bytes in its 64-byte FIFO, starts a \emph{Transceive} command, and reads the
card's answer back from the FIFO. Reading a card's UID follows ISO/IEC 14443-3 type A:
\begin{enumerate}
  \item \textbf{REQA} (or \textbf{WUPA} for a card that was put to sleep): ``is a card there?''.
        A card answers with its 2-byte ATQA.
  \item \textbf{ANTICOLL}: the card sends 4 UID bytes and their XOR check byte (BCC), which the
        driver verifies.
  \item \textbf{SELECT} with a CRC\_A: the card confirms and answers SAK, which says whether more
        UID bytes follow (7-byte UIDs take two such levels).
  \item \textbf{HLTA}: the card goes to sleep, so a card left on the reader is not read again and
        again.
\end{enumerate}
\file{components/nodes/rfid.c} owns SPI bus~3 and both readers. Its task polls each reader every
100\,ms. A card counts as a \emph{new} presentation when that reader has not seen it during the last
1.5\,s, so a card held on the reader opens the door once, and presenting it again later is a new
scan. A 4-byte UID is used as the 32-bit card number first byte first (Blue card
\code{01:02:03:04} $\rightarrow$ \code{01020304}); a 7-byte UID keeps its last 4 bytes. At boot
the task resets both readers and checks each chip answers; a reader that does not is logged as
\code{ALERT code=RFID_FAIL} and retried every 5\,s, and \code{status} shows the readers' health.

\subsection{Pure decision logic (\code{core_logic})}
The rules that matter most are plain C functions with no hardware access: they take a struct of
facts and return a decision. This keeps the rules short, readable and testable on a PC
(\cref{sec:testing}). The access policy is the clearest example:

\begin{lstlisting}
access_reason_t access_decide(const access_request_t *req)
{
    if (!req->known)   return ACCESS_UNKNOWN_CARD;
    if (!req->enabled) return ACCESS_DISABLED;
    if ((req->side == SIDE_IN && req->inside) || (req->side == SIDE_OUT && !req->inside))
        return ACCESS_ANTI_PASSBACK;
    bool privileged = role_is_privileged(req->role);      /* SECURITY or ADMIN */
    if (req->lockdown)      return privileged ? ACCESS_OK_PRIVILEGED : ACCESS_LOCKDOWN;
    if (req->schedule_open) return ACCESS_OK;
    if (privileged)         return ACCESS_OK_PRIVILEGED;
    if (req->side == SIDE_OUT && req->escort_active) return ACCESS_OK_ESCORT;
    return ACCESS_OUTSIDE_HOURS;
}
\end{lstlisting}
(Shown compacted; the real function in \file{access_policy.c} has the same logic.)
\code{room_logic_step()} does the same for lighting (\cref{sec:op-lighting}) and
\code{auth_fsm} for the admin login (\cref{sec:op-auth}).

\subsection{Shared state services}
Each module in \file{components/services/} owns one piece of shared state, guards it with its own
mutex, hands out copies to readers, and saves it to NVS (the ESP32's key-value flash store) so
it survives a reboot:
\begin{itemize}
  \item \code{users}: the card registry (UID, name, role, enabled, salted SHA-256 PIN hash). It is
        versioned, so a changed layout re-seeds the four demo cards.
  \item \code{config}: the access schedule and lockdown flag. A change is saved to NVS
        \emph{before} it becomes live, so a power cut never leaves the two disagreeing.
  \item \code{occupancy}: the ledger of who is inside, saved on every change.
  \item \code{auth}: wraps the pure \code{auth_fsm} and writes every login result to the log.
  \item \code{clock}: sets the system time from the DS1307 at boot (time zone EAT, UTC+3).
\end{itemize}

\subsection{Interrupts}
Interrupt handlers do the minimum: they push the pin number to a queue or notify a task, and
return. Debouncing, scanning and decisions happen in tasks. This keeps interrupts short and makes it
safe to call any FreeRTOS or logging function from the code that handles the press.

\subsection{Coding principles}
\begin{itemize}
  \item \textbf{Single source of truth.} Pins are only defined in \file{board_pins.h}, task settings
        only in \file{task_config.h}, tunable behaviour only in \file{Kconfig.projbuild}, and each
        enum and its string names only once, using X-macros in \file{cocis_types.h}.
  \item \textbf{One way to do each thing.} All modules log through
        \code{cocis_log(type, "k=v ...")}. One PWM layer drives servos, lamps and the buzzer. One
        button driver (interrupt, queue, debounce) serves any push button. Every card scan, from
        either reader or the CLI's \code{tap}, goes through the same \code{rfid_inject()}. One NVS
        blob helper is used by all persistent modules.
  \item \textbf{Separation of concerns.} Drivers know hardware, services know state, nodes combine
        them, and \code{core_logic} knows only the rules.
\end{itemize}
